% Ngân hàng câu hỏi Toán 10 (KNTT) -- Chương 7, Bài 19, Dạng 45
%% Dạng: Viết phương trình đường thẳng
%% 100 câu: NB 30 -- TH 40 -- VD 20 -- VDC 10 ; mỗi mức: trắc nghiệm : đúng/sai : trả lời ngắn = 3 : 1 : 1
%% Phương án ĐẦU TIÊN trong \cauTN luôn là đáp án đúng (gói trondethi tự đảo).
%% Phần II: \y{nội dung}{d|s}. Phần III: đáp án đặt ở đối số cuối.

\cauTN[id=T10-C7-B19-D45-NB-001,muc=NB,loigiai={$(x-6)+3(y-6)=0\Leftrightarrow x+3y-24=0$.}]{Phương trình tổng quát của đường thẳng đi qua $M(6;6)$ và có vectơ pháp tuyến $\vec n=(1;3)$ là}
  {$x+3y-24=0$.}
  {$x+3y+24=0$.}{$3x+y-24=0$.}{$x-3y-24=0$.}

\cauTN[id=T10-C7-B19-D45-NB-002,muc=NB,loigiai={$\overrightarrow{AB}=(7;-1)$, vectơ pháp tuyến $(1;7)$, đi qua $A$: $x+7y-34=0$.}]{Phương trình tổng quát của đường thẳng $AB$ với $A(-1;5)$, $B(6;4)$ là}
  {$x+7y-34=0$.}
  {$x+7y+34=0$.}{$7x+y-34=0$.}{$x-7y-34=0$.}

\cauTN[id=T10-C7-B19-D45-NB-003,muc=NB,loigiai={$2(x-4)+4(y-5)=0\Leftrightarrow x+2y-14=0$.}]{Phương trình tổng quát của đường thẳng đi qua $M(4;5)$ và có vectơ pháp tuyến $\vec n=(2;4)$ là}
  {$x+2y-14=0$.}
  {$x+2y+14=0$.}{$2x+y-14=0$.}{$x-2y-14=0$.}

\cauTN[id=T10-C7-B19-D45-NB-004,muc=NB,loigiai={$x=x_0+at$, $y=y_0+bt$.}]{Phương trình tham số của đường thẳng đi qua $M(2;-2)$ và có vectơ chỉ phương $\vec u=(-4;-1)$ là}
  {$\left\{\begin{array}{l}x=-4t+2\\y=-t-2\end{array}\right.$.}
  {$\left\{\begin{array}{l}x=2t-4\\y=-2t-1\end{array}\right.$.}{$\left\{\begin{array}{l}x=-t+2\\y=4t-2\end{array}\right.$.}{$\left\{\begin{array}{l}x=4t+2\\y=-t-2\end{array}\right.$.}

\cauTN[id=T10-C7-B19-D45-NB-005,muc=NB,loigiai={$\overrightarrow{AB}=(7;5)$, vectơ pháp tuyến $(5;-7)$, đi qua $A$: $5x-7y-22=0$.}]{Phương trình tổng quát của đường thẳng $AB$ với $A(-4;-6)$, $B(3;-1)$ là}
  {$5x-7y-22=0$.}
  {$5x-7y+22=0$.}{$7x-5y+22=0$.}{$5x+7y-22=0$.}

\cauTN[id=T10-C7-B19-D45-NB-006,muc=NB,loigiai={$-5(x-5)+4(y-1)=0\Leftrightarrow 5x-4y-21=0$.}]{Phương trình tổng quát của đường thẳng đi qua $M(5;1)$ và có vectơ pháp tuyến $\vec n=(-5;4)$ là}
  {$5x-4y-21=0$.}
  {$5x-4y+21=0$.}{$4x-5y+21=0$.}{$5x+4y-21=0$.}

\cauTN[id=T10-C7-B19-D45-NB-007,muc=NB,loigiai={$\overrightarrow{AB}=(3;-11)$, vectơ pháp tuyến $(11;3)$, đi qua $A$: $11x+3y+29=0$.}]{Phương trình tổng quát của đường thẳng $AB$ với $A(-4;5)$, $B(-1;-6)$ là}
  {$11x+3y+29=0$.}
  {$11x+3y-29=0$.}{$3x+11y+29=0$.}{$11x-3y+29=0$.}

\cauTN[id=T10-C7-B19-D45-NB-008,muc=NB,loigiai={$-5(x-2)-3(y-1)=0\Leftrightarrow 5x+3y-13=0$.}]{Phương trình tổng quát của đường thẳng đi qua $M(2;1)$ và có vectơ pháp tuyến $\vec n=(-5;-3)$ là}
  {$5x+3y-13=0$.}
  {$5x+3y+13=0$.}{$3x+5y-13=0$.}{$5x-3y-13=0$.}

\cauTN[id=T10-C7-B19-D45-NB-009,muc=NB,loigiai={$x=x_0+at$, $y=y_0+bt$.}]{Phương trình tham số của đường thẳng đi qua $M(-6;3)$ và có vectơ chỉ phương $\vec u=(-2;-5)$ là}
  {$\left\{\begin{array}{l}x=-2t-6\\y=-5t+3\end{array}\right.$.}
  {$\left\{\begin{array}{l}x=-6t-2\\y=3t-5\end{array}\right.$.}{$\left\{\begin{array}{l}x=-5t-6\\y=2t+3\end{array}\right.$.}{$\left\{\begin{array}{l}x=2t-6\\y=-5t+3\end{array}\right.$.}

\cauTN[id=T10-C7-B19-D45-NB-010,muc=NB,loigiai={$\overrightarrow{AB}=(2;-5)$, vectơ pháp tuyến $(5;2)$, đi qua $A$: $5x+2y-15=0$.}]{Phương trình tổng quát của đường thẳng $AB$ với $A(1;5)$, $B(3;0)$ là}
  {$5x+2y-15=0$.}
  {$5x+2y+15=0$.}{$2x+5y-15=0$.}{$5x-2y-15=0$.}

\cauTN[id=T10-C7-B19-D45-NB-011,muc=NB,loigiai={$3(x-1)+5(y-5)=0\Leftrightarrow 3x+5y-28=0$.}]{Phương trình tổng quát của đường thẳng đi qua $M(1;5)$ và có vectơ pháp tuyến $\vec n=(3;5)$ là}
  {$3x+5y-28=0$.}
  {$3x+5y+28=0$.}{$5x+3y-28=0$.}{$3x-5y-28=0$.}

\cauTN[id=T10-C7-B19-D45-NB-012,muc=NB,loigiai={$\overrightarrow{AB}=(-6;-8)$, vectơ pháp tuyến $(4;-3)$, đi qua $A$: $4x-3y+12=0$.}]{Phương trình tổng quát của đường thẳng $AB$ với $A(0;4)$, $B(-6;-4)$ là}
  {$4x-3y+12=0$.}
  {$4x-3y-12=0$.}{$3x-4y-12=0$.}{$4x+3y+12=0$.}

\cauTN[id=T10-C7-B19-D45-NB-013,muc=NB,loigiai={$1(x)-2(y+2)=0\Leftrightarrow x-2y-4=0$.}]{Phương trình tổng quát của đường thẳng đi qua $M(0;-2)$ và có vectơ pháp tuyến $\vec n=(1;-2)$ là}
  {$x-2y-4=0$.}
  {$x-2y+4=0$.}{$2x-y+4=0$.}{$x+2y-4=0$.}

\cauTN[id=T10-C7-B19-D45-NB-014,muc=NB,loigiai={$x=x_0+at$, $y=y_0+bt$.}]{Phương trình tham số của đường thẳng đi qua $M(-2;-3)$ và có vectơ chỉ phương $\vec u=(4;-1)$ là}
  {$\left\{\begin{array}{l}x=4t-2\\y=-t-3\end{array}\right.$.}
  {$\left\{\begin{array}{l}x=-2t+4\\y=-3t-1\end{array}\right.$.}{$\left\{\begin{array}{l}x=-t-2\\y=-4t-3\end{array}\right.$.}{$\left\{\begin{array}{l}x=-4t-2\\y=-t-3\end{array}\right.$.}

\cauTN[id=T10-C7-B19-D45-NB-015,muc=NB,loigiai={$\overrightarrow{AB}=(-1;-9)$, vectơ pháp tuyến $(9;-1)$, đi qua $A$: $9x-y-24=0$.}]{Phương trình tổng quát của đường thẳng $AB$ với $A(3;3)$, $B(2;-6)$ là}
  {$9x-y-24=0$.}
  {$9x-y+24=0$.}{$x-9y+24=0$.}{$9x+y-24=0$.}

\cauTN[id=T10-C7-B19-D45-NB-016,muc=NB,loigiai={$(x-6)+5(y)=0\Leftrightarrow x+5y-6=0$.}]{Phương trình tổng quát của đường thẳng đi qua $M(6;0)$ và có vectơ pháp tuyến $\vec n=(1;5)$ là}
  {$x+5y-6=0$.}
  {$x+5y+6=0$.}{$5x+y-6=0$.}{$x-5y-6=0$.}

\cauTN[id=T10-C7-B19-D45-NB-017,muc=NB,loigiai={$x=x_0+at$, $y=y_0+bt$.}]{Phương trình tham số của đường thẳng đi qua $M(-3;-5)$ và có vectơ chỉ phương $\vec u=(-2;3)$ là}
  {$\left\{\begin{array}{l}x=-2t-3\\y=3t-5\end{array}\right.$.}
  {$\left\{\begin{array}{l}x=-3t-2\\y=-5t+3\end{array}\right.$.}{$\left\{\begin{array}{l}x=3t-3\\y=2t-5\end{array}\right.$.}{$\left\{\begin{array}{l}x=2t-3\\y=3t-5\end{array}\right.$.}

\cauTN[id=T10-C7-B19-D45-NB-018,muc=NB,loigiai={$\overrightarrow{AB}=(3;2)$, vectơ pháp tuyến $(2;-3)$, đi qua $A$: $2x-3y+3=0$.}]{Phương trình tổng quát của đường thẳng $AB$ với $A(0;1)$, $B(3;3)$ là}
  {$2x-3y+3=0$.}
  {$2x-3y-3=0$.}{$3x-2y-3=0$.}{$2x+3y+3=0$.}

\cauDS[id=T10-C7-B19-D45-NB-019,muc=NB,loigiai={$AB\colon 2x+3y-6=0$; điểm đối xứng của $A$ qua $B$ thuộc $AB$. a) Đúng. b) Sai. c) Đúng. d) Sai.}]{Cho hai điểm $A(0;2)$, $B(6;-2)$.}{%
  \y{$\overrightarrow{AB}=(6;-4)$ là một vectơ chỉ phương của đường thẳng $AB$.}{d}
  \y{Đường thẳng $AB$ có phương trình $2x-3y-6=0$.}{s}
  \y{Điểm $(12;-6)$ thuộc đường thẳng $AB$.}{d}
  \y{Đường thẳng $AB$ đi qua gốc toạ độ.}{s}}

\cauDS[id=T10-C7-B19-D45-NB-020,muc=NB,loigiai={$AB\colon x-y=0$; điểm đối xứng của $A$ qua $B$ thuộc $AB$. a) Đúng. b) Sai. c) Đúng. d) Đúng.}]{Cho hai điểm $A(5;5)$, $B(6;6)$.}{%
  \y{$\overrightarrow{AB}=(1;1)$ là một vectơ chỉ phương của đường thẳng $AB$.}{d}
  \y{Đường thẳng $AB$ có phương trình $x+y=0$.}{s}
  \y{Điểm $(7;7)$ thuộc đường thẳng $AB$.}{d}
  \y{Đường thẳng $AB$ đi qua gốc toạ độ.}{d}}

\cauDS[id=T10-C7-B19-D45-NB-021,muc=NB,loigiai={$AB\colon 2x-3y-11=0$; điểm đối xứng của $A$ qua $B$ thuộc $AB$. a) Đúng. b) Sai. c) Đúng. d) Sai.}]{Cho hai điểm $A(1;-3)$, $B(4;-1)$.}{%
  \y{$\overrightarrow{AB}=(3;2)$ là một vectơ chỉ phương của đường thẳng $AB$.}{d}
  \y{Đường thẳng $AB$ có phương trình $2x+3y-11=0$.}{s}
  \y{Điểm $(7;1)$ thuộc đường thẳng $AB$.}{d}
  \y{Đường thẳng $AB$ đi qua gốc toạ độ.}{s}}

\cauDS[id=T10-C7-B19-D45-NB-022,muc=NB,loigiai={$AB\colon 2x+3y+9=0$; điểm đối xứng của $A$ qua $B$ thuộc $AB$. a) Đúng. b) Đúng. c) Đúng. d) Sai.}]{Cho hai điểm $A(3;-5)$, $B(0;-3)$.}{%
  \y{$\overrightarrow{AB}=(-3;2)$ là một vectơ chỉ phương của đường thẳng $AB$.}{d}
  \y{Đường thẳng $AB$ có phương trình $2x+3y+9=0$.}{d}
  \y{Điểm $(-3;-1)$ thuộc đường thẳng $AB$.}{d}
  \y{Đường thẳng $AB$ đi qua gốc toạ độ.}{s}}

\cauDS[id=T10-C7-B19-D45-NB-023,muc=NB,loigiai={$AB\colon x+6=0$; điểm đối xứng của $A$ qua $B$ thuộc $AB$. a) Đúng. b) Đúng. c) Đúng. d) Sai.}]{Cho hai điểm $A(-6;-6)$, $B(-6;-1)$.}{%
  \y{$\overrightarrow{AB}=(0;5)$ là một vectơ chỉ phương của đường thẳng $AB$.}{d}
  \y{Đường thẳng $AB$ có phương trình $x+6=0$.}{d}
  \y{Điểm $(-6;4)$ thuộc đường thẳng $AB$.}{d}
  \y{Đường thẳng $AB$ đi qua gốc toạ độ.}{s}}

\cauDS[id=T10-C7-B19-D45-NB-024,muc=NB,loigiai={$AB\colon x+2y-1=0$; điểm đối xứng của $A$ qua $B$ thuộc $AB$. a) Đúng. b) Sai. c) Đúng. d) Sai.}]{Cho hai điểm $A(-5;3)$, $B(-3;2)$.}{%
  \y{$\overrightarrow{AB}=(2;-1)$ là một vectơ chỉ phương của đường thẳng $AB$.}{d}
  \y{Đường thẳng $AB$ có phương trình $x-2y-1=0$.}{s}
  \y{Điểm $(-1;1)$ thuộc đường thẳng $AB$.}{d}
  \y{Đường thẳng $AB$ đi qua gốc toạ độ.}{s}}

\cauTLN[id=T10-C7-B19-D45-NB-025,muc=NB,loigiai={Phương trình $2x-3y-15=0$, chia cho $2$: $b=-\dfrac{3}{2}$, $c=-\dfrac{15}{2}$.}]{Đường thẳng đi qua $A(3;-3)$ và $B(6;-1)$ có phương trình dạng $x+by+c=0$. Tính $b+c$.}{-9}

\cauTLN[id=T10-C7-B19-D45-NB-026,muc=NB,loigiai={Phương trình $x-8y+14=0$, chia cho $1$: $b=-8$, $c=14$.}]{Đường thẳng đi qua $A(-6;1)$ và $B(2;2)$ có phương trình dạng $x+by+c=0$. Tính $b+c$.}{6}

\cauTLN[id=T10-C7-B19-D45-NB-027,muc=NB,loigiai={Phương trình $x-y+1=0$, chia cho $1$: $b=-1$, $c=1$.}]{Đường thẳng đi qua $A(1;2)$ và $B(-6;-5)$ có phương trình dạng $x+by+c=0$. Tính $b+c$.}{0}

\cauTLN[id=T10-C7-B19-D45-NB-028,muc=NB,loigiai={Phương trình $x+y=0$, chia cho $1$: $b=1$, $c=0$.}]{Đường thẳng đi qua $A(1;-1)$ và $B(3;-3)$ có phương trình dạng $x+by+c=0$. Tính $b+c$.}{1}

\cauTLN[id=T10-C7-B19-D45-NB-029,muc=NB,loigiai={Phương trình $x-6y-1=0$, chia cho $1$: $b=-6$, $c=-1$.}]{Đường thẳng đi qua $A(1;0)$ và $B(-5;-1)$ có phương trình dạng $x+by+c=0$. Tính $b+c$.}{-7}

\cauTLN[id=T10-C7-B19-D45-NB-030,muc=NB,loigiai={Phương trình $x+3=0$, chia cho $1$: $b=0$, $c=3$.}]{Đường thẳng đi qua $A(-3;6)$ và $B(-3;-2)$ có phương trình dạng $x+by+c=0$. Tính $b+c$.}{3}

\cauTN[id=T10-C7-B19-D45-TH-031,muc=TH,loigiai={Cùng vectơ pháp tuyến $(2;-3)$; đi qua $M$: $2x-3y+9=0$.}]{Đường thẳng đi qua $M(-6;-1)$ và song song với đường thẳng $2x-3y-1=0$ có phương trình là}
  {$2x-3y+9=0$.}
  {$2x-3y-9=0$.}{$3x-2y-9=0$.}{$2x+3y+9=0$.}

\cauTN[id=T10-C7-B19-D45-TH-032,muc=TH,loigiai={Đi qua trung điểm $I(\dfrac{7}{2};-\dfrac{3}{2})$, vectơ pháp tuyến $\overrightarrow{AB}=(-5;-3)$: $5x+3y-13=0$.}]{Phương trình đường trung trực của đoạn thẳng $AB$ với $A(6;0)$, $B(1;-3)$ là}
  {$5x+3y-13=0$.}
  {$5x+3y+13=0$.}{$3x+5y-13=0$.}{$5x-3y-13=0$.}

\cauTN[id=T10-C7-B19-D45-TH-033,muc=TH,loigiai={Phương trình đoạn chắn: $\dfrac xa+\dfrac yb=1$.}]{Đường thẳng cắt trục $Ox$ tại $A(6;0)$ và cắt trục $Oy$ tại $B(0;-3)$ có phương trình là}
  {$\dfrac{x}{6}-\dfrac{y}{3}=1$.}
  {$-\dfrac{x}{3}+\dfrac{y}{6}=1$.}{$\dfrac{x}{6}-\dfrac{y}{3}=0$.}{$6x-3y-1=0$.}

\cauTN[id=T10-C7-B19-D45-TH-034,muc=TH,loigiai={Cùng vectơ pháp tuyến $(1;-4)$; đi qua $M$: $x-4y-1=0$.}]{Đường thẳng đi qua $M(1;0)$ và song song với đường thẳng $x-4y+3=0$ có phương trình là}
  {$x-4y-1=0$.}
  {$x-4y+1=0$.}{$4x-y+1=0$.}{$x+4y-1=0$.}

\cauTN[id=T10-C7-B19-D45-TH-035,muc=TH,loigiai={Đi qua trung điểm $I(-4;-\dfrac{7}{2})$, vectơ pháp tuyến $\overrightarrow{AB}=(2;1)$: $4x+2y+23=0$.}]{Phương trình đường trung trực của đoạn thẳng $AB$ với $A(-5;-4)$, $B(-3;-3)$ là}
  {$4x+2y+23=0$.}
  {$4x+2y-23=0$.}{$2x+4y+23=0$.}{$4x-2y+23=0$.}

\cauTN[id=T10-C7-B19-D45-TH-036,muc=TH,loigiai={Phương trình đoạn chắn: $\dfrac xa+\dfrac yb=1$.}]{Đường thẳng cắt trục $Ox$ tại $A(5;0)$ và cắt trục $Oy$ tại $B(0;1)$ có phương trình là}
  {$\dfrac{x}{5}+\dfrac{y}{1}=1$.}
  {$\dfrac{x}{1}+\dfrac{y}{5}=1$.}{$\dfrac{x}{5}+\dfrac{y}{1}=0$.}{$5x+y-1=0$.}

\cauTN[id=T10-C7-B19-D45-TH-037,muc=TH,loigiai={Cùng vectơ pháp tuyến $(3;1)$; đi qua $M$: $3x+y-1=0$.}]{Đường thẳng đi qua $M(0;1)$ và song song với đường thẳng $3x+y+4=0$ có phương trình là}
  {$3x+y-1=0$.}
  {$3x+y+1=0$.}{$x+3y-1=0$.}{$3x-y-1=0$.}

\cauTN[id=T10-C7-B19-D45-TH-038,muc=TH,loigiai={Đi qua trung điểm $I(4;0)$, vectơ pháp tuyến $\overrightarrow{AB}=(-2;-12)$: $x+6y-4=0$.}]{Phương trình đường trung trực của đoạn thẳng $AB$ với $A(5;6)$, $B(3;-6)$ là}
  {$x+6y-4=0$.}
  {$x+6y+4=0$.}{$6x+y-4=0$.}{$x-6y-4=0$.}

\cauTN[id=T10-C7-B19-D45-TH-039,muc=TH,loigiai={Phương trình đoạn chắn: $\dfrac xa+\dfrac yb=1$.}]{Đường thẳng cắt trục $Ox$ tại $A(3;0)$ và cắt trục $Oy$ tại $B(0;-1)$ có phương trình là}
  {$\dfrac{x}{3}-\dfrac{y}{1}=1$.}
  {$-\dfrac{x}{1}+\dfrac{y}{3}=1$.}{$\dfrac{x}{3}-\dfrac{y}{1}=0$.}{$3x-y-1=0$.}

\cauTN[id=T10-C7-B19-D45-TH-040,muc=TH,loigiai={Cùng vectơ pháp tuyến $(3;-5)$; đi qua $M$: $3x-5y-34=0$.}]{Đường thẳng đi qua $M(3;-5)$ và song song với đường thẳng $3x-5y+2=0$ có phương trình là}
  {$3x-5y-34=0$.}
  {$3x-5y+34=0$.}{$5x-3y+34=0$.}{$3x+5y-34=0$.}

\cauTN[id=T10-C7-B19-D45-TH-041,muc=TH,loigiai={Đi qua trung điểm $I(\dfrac{9}{2};\dfrac{3}{2})$, vectơ pháp tuyến $\overrightarrow{AB}=(1;3)$: $x+3y-9=0$.}]{Phương trình đường trung trực của đoạn thẳng $AB$ với $A(4;0)$, $B(5;3)$ là}
  {$x+3y-9=0$.}
  {$x+3y+9=0$.}{$3x+y-9=0$.}{$x-3y-9=0$.}

\cauTN[id=T10-C7-B19-D45-TH-042,muc=TH,loigiai={Phương trình đoạn chắn: $\dfrac xa+\dfrac yb=1$.}]{Đường thẳng cắt trục $Ox$ tại $A(-5;0)$ và cắt trục $Oy$ tại $B(0;-5)$ có phương trình là}
  {$-\dfrac{x}{5}-\dfrac{y}{5}=1$.}
  {$-\dfrac{x}{5}+\dfrac{y}{5}=1$.}{$-\dfrac{x}{5}-\dfrac{y}{5}=0$.}{$5x+5y+1=0$.}

\cauTN[id=T10-C7-B19-D45-TH-043,muc=TH,loigiai={Cùng vectơ pháp tuyến $(3;2)$; đi qua $M$: $3x+2y-7=0$.}]{Đường thẳng đi qua $M(1;2)$ và song song với đường thẳng $3x+2y+2=0$ có phương trình là}
  {$3x+2y-7=0$.}
  {$3x+2y+7=0$.}{$2x+3y-7=0$.}{$3x-2y-7=0$.}

\cauTN[id=T10-C7-B19-D45-TH-044,muc=TH,loigiai={Đi qua trung điểm $I(-\dfrac{9}{2};-\dfrac{1}{2})$, vectơ pháp tuyến $\overrightarrow{AB}=(1;3)$: $x+3y+6=0$.}]{Phương trình đường trung trực của đoạn thẳng $AB$ với $A(-5;-2)$, $B(-4;1)$ là}
  {$x+3y+6=0$.}
  {$x+3y-6=0$.}{$3x+y+6=0$.}{$x-3y+6=0$.}

\cauTN[id=T10-C7-B19-D45-TH-045,muc=TH,loigiai={Phương trình đoạn chắn: $\dfrac xa+\dfrac yb=1$.}]{Đường thẳng cắt trục $Ox$ tại $A(3;0)$ và cắt trục $Oy$ tại $B(0;-3)$ có phương trình là}
  {$\dfrac{x}{3}-\dfrac{y}{3}=1$.}
  {$-\dfrac{x}{3}+\dfrac{y}{3}=1$.}{$\dfrac{x}{3}-\dfrac{y}{3}=0$.}{$3x-3y-1=0$.}

\cauTN[id=T10-C7-B19-D45-TH-046,muc=TH,loigiai={Vectơ pháp tuyến là vectơ chỉ phương $(3;-5)$ của đường đã cho; đi qua $M$: $3x-5y-33=0$.}]{Đường thẳng đi qua $M(1;-6)$ và vuông góc với đường thẳng $5x+3y+6=0$ có phương trình là}
  {$3x-5y-33=0$.}
  {$3x-5y+33=0$.}{$5x-3y+33=0$.}{$3x+5y-33=0$.}

\cauTN[id=T10-C7-B19-D45-TH-047,muc=TH,loigiai={Đi qua trung điểm $I(\dfrac{1}{2};\dfrac{11}{2})$, vectơ pháp tuyến $\overrightarrow{AB}=(9;-1)$: $9x-y+1=0$.}]{Phương trình đường trung trực của đoạn thẳng $AB$ với $A(-4;6)$, $B(5;5)$ là}
  {$9x-y+1=0$.}
  {$9x-y-1=0$.}{$x-9y-1=0$.}{$9x+y+1=0$.}

\cauTN[id=T10-C7-B19-D45-TH-048,muc=TH,loigiai={Vectơ pháp tuyến là vectơ chỉ phương $(-1;-1)$ của đường đã cho; đi qua $M$: $x+y-8=0$.}]{Đường thẳng đi qua $M(5;3)$ và vuông góc với đường thẳng $x-y-2=0$ có phương trình là}
  {$x+y-8=0$.}
  {$x+y+8=0$.}{$x-y-8=0$.}{$x+y-6=0$.}

\cauTN[id=T10-C7-B19-D45-TH-049,muc=TH,loigiai={Đi qua trung điểm $I(-\dfrac{3}{2};-\dfrac{7}{2})$, vectơ pháp tuyến $\overrightarrow{AB}=(5;5)$: $x+y+5=0$.}]{Phương trình đường trung trực của đoạn thẳng $AB$ với $A(-4;-6)$, $B(1;-1)$ là}
  {$x+y+5=0$.}
  {$x+y-5=0$.}{$x-y+5=0$.}{$x+y+7=0$.}

\cauTN[id=T10-C7-B19-D45-TH-050,muc=TH,loigiai={Phương trình đoạn chắn: $\dfrac xa+\dfrac yb=1$.}]{Đường thẳng cắt trục $Ox$ tại $A(5;0)$ và cắt trục $Oy$ tại $B(0;-4)$ có phương trình là}
  {$\dfrac{x}{5}-\dfrac{y}{4}=1$.}
  {$-\dfrac{x}{4}+\dfrac{y}{5}=1$.}{$\dfrac{x}{5}-\dfrac{y}{4}=0$.}{$5x-4y-1=0$.}

\cauTN[id=T10-C7-B19-D45-TH-051,muc=TH,loigiai={Vectơ pháp tuyến là vectơ chỉ phương $(-4;-3)$ của đường đã cho; đi qua $M$: $4x+3y-13=0$.}]{Đường thẳng đi qua $M(1;3)$ và vuông góc với đường thẳng $3x-4y+1=0$ có phương trình là}
  {$4x+3y-13=0$.}
  {$4x+3y+13=0$.}{$3x+4y-13=0$.}{$4x-3y-13=0$.}

\cauTN[id=T10-C7-B19-D45-TH-052,muc=TH,loigiai={Đi qua trung điểm $I(\dfrac{7}{2};-\dfrac{5}{2})$, vectơ pháp tuyến $\overrightarrow{AB}=(-3;-7)$: $3x+7y+7=0$.}]{Phương trình đường trung trực của đoạn thẳng $AB$ với $A(5;1)$, $B(2;-6)$ là}
  {$3x+7y+7=0$.}
  {$3x+7y-7=0$.}{$7x+3y+7=0$.}{$3x-7y+7=0$.}

\cauTN[id=T10-C7-B19-D45-TH-053,muc=TH,loigiai={Phương trình đoạn chắn: $\dfrac xa+\dfrac yb=1$.}]{Đường thẳng cắt trục $Ox$ tại $A(-4;0)$ và cắt trục $Oy$ tại $B(0;-6)$ có phương trình là}
  {$-\dfrac{x}{4}-\dfrac{y}{6}=1$.}
  {$-\dfrac{x}{6}-\dfrac{y}{4}=1$.}{$-\dfrac{x}{4}-\dfrac{y}{6}=0$.}{$4x+6y+1=0$.}

\cauTN[id=T10-C7-B19-D45-TH-054,muc=TH,loigiai={Cùng vectơ pháp tuyến $(4;5)$; đi qua $M$: $4x+5y+14=0$.}]{Đường thẳng đi qua $M(-1;-2)$ và song song với đường thẳng $4x+5y+3=0$ có phương trình là}
  {$4x+5y+14=0$.}
  {$4x+5y-14=0$.}{$5x+4y+14=0$.}{$4x-5y+14=0$.}

\cauDS[id=T10-C7-B19-D45-TH-055,muc=TH,loigiai={$BC\colon 8x-y+46=0$; đường cao từ $A$: $x+8y+52=0$. a) Đúng. b) Đúng. c) Đúng. d) Sai.}]{Cho tam giác $ABC$ với $A(-4;-6)$, $B(-6;-2)$, $C(-5;6)$.}{%
  \y{Đường thẳng $BC$ có phương trình $8x-y+46=0$.}{d}
  \y{Đường cao kẻ từ $A$ có vectơ pháp tuyến $\overrightarrow{BC}=(1;8)$.}{d}
  \y{Đường cao kẻ từ $A$ có phương trình $x+8y+52=0$.}{d}
  \y{Đường cao kẻ từ $A$ song song với $BC$.}{s}}

\cauDS[id=T10-C7-B19-D45-TH-056,muc=TH,loigiai={$AB\colon x+5=0$; điểm đối xứng của $A$ qua $B$ thuộc $AB$. a) Đúng. b) Sai. c) Đúng. d) Sai.}]{Cho hai điểm $A(-5;4)$, $B(-5;5)$.}{%
  \y{$\overrightarrow{AB}=(0;1)$ là một vectơ chỉ phương của đường thẳng $AB$.}{d}
  \y{Đường thẳng $AB$ có phương trình $x+y+5=0$.}{s}
  \y{Điểm $(-5;6)$ thuộc đường thẳng $AB$.}{d}
  \y{Đường thẳng $AB$ đi qua gốc toạ độ.}{s}}

\cauDS[id=T10-C7-B19-D45-TH-057,muc=TH,loigiai={$BC\colon 9x+4y-6=0$; đường cao từ $A$: $4x-9y-27=0$. a) Đúng. b) Đúng. c) Đúng. d) Sai.}]{Cho tam giác $ABC$ với $A(0;-3)$, $B(2;-3)$, $C(-2;6)$.}{%
  \y{Đường thẳng $BC$ có phương trình $9x+4y-6=0$.}{d}
  \y{Đường cao kẻ từ $A$ có vectơ pháp tuyến $\overrightarrow{BC}=(-4;9)$.}{d}
  \y{Đường cao kẻ từ $A$ có phương trình $4x-9y-27=0$.}{d}
  \y{Đường cao kẻ từ $A$ song song với $BC$.}{s}}

\cauDS[id=T10-C7-B19-D45-TH-058,muc=TH,loigiai={$AB\colon 5x+11y-3=0$; điểm đối xứng của $A$ qua $B$ thuộc $AB$. a) Đúng. b) Sai. c) Đúng. d) Sai.}]{Cho hai điểm $A(-6;3)$, $B(5;-2)$.}{%
  \y{$\overrightarrow{AB}=(11;-5)$ là một vectơ chỉ phương của đường thẳng $AB$.}{d}
  \y{Đường thẳng $AB$ có phương trình $5x-11y-3=0$.}{s}
  \y{Điểm $(16;-7)$ thuộc đường thẳng $AB$.}{d}
  \y{Đường thẳng $AB$ đi qua gốc toạ độ.}{s}}

\cauDS[id=T10-C7-B19-D45-TH-059,muc=TH,loigiai={$BC\colon 3x+y+9=0$; đường cao từ $A$: $x-3y+21=0$. a) Đúng. b) Đúng. c) Đúng. d) Sai.}]{Cho tam giác $ABC$ với $A(-6;5)$, $B(-5;6)$, $C(-4;3)$.}{%
  \y{Đường thẳng $BC$ có phương trình $3x+y+9=0$.}{d}
  \y{Đường cao kẻ từ $A$ có vectơ pháp tuyến $\overrightarrow{BC}=(1;-3)$.}{d}
  \y{Đường cao kẻ từ $A$ có phương trình $x-3y+21=0$.}{d}
  \y{Đường cao kẻ từ $A$ song song với $BC$.}{s}}

\cauDS[id=T10-C7-B19-D45-TH-060,muc=TH,loigiai={$AB\colon x+3y+16=0$; điểm đối xứng của $A$ qua $B$ thuộc $AB$. a) Đúng. b) Đúng. c) Đúng. d) Sai.}]{Cho hai điểm $A(2;-6)$, $B(-4;-4)$.}{%
  \y{$\overrightarrow{AB}=(-6;2)$ là một vectơ chỉ phương của đường thẳng $AB$.}{d}
  \y{Đường thẳng $AB$ có phương trình $x+3y+16=0$.}{d}
  \y{Điểm $(-10;-2)$ thuộc đường thẳng $AB$.}{d}
  \y{Đường thẳng $AB$ đi qua gốc toạ độ.}{s}}

\cauDS[id=T10-C7-B19-D45-TH-061,muc=TH,loigiai={$BC\colon 2x+5y+20=0$; đường cao từ $A$: $5x-2y+27=0$. a) Đúng. b) Đúng. c) Sai. d) Sai.}]{Cho tam giác $ABC$ với $A(-3;6)$, $B(5;-6)$, $C(-5;-2)$.}{%
  \y{Đường thẳng $BC$ có phương trình $2x+5y+20=0$.}{d}
  \y{Đường cao kẻ từ $A$ có vectơ pháp tuyến $\overrightarrow{BC}=(-10;4)$.}{d}
  \y{Đường cao kẻ từ $A$ có phương trình $5x-2y+28=0$.}{s}
  \y{Đường cao kẻ từ $A$ song song với $BC$.}{s}}

\cauDS[id=T10-C7-B19-D45-TH-062,muc=TH,loigiai={$AB\colon 5x-8y+10=0$; điểm đối xứng của $A$ qua $B$ thuộc $AB$. a) Đúng. b) Đúng. c) Đúng. d) Sai.}]{Cho hai điểm $A(-2;0)$, $B(6;5)$.}{%
  \y{$\overrightarrow{AB}=(8;5)$ là một vectơ chỉ phương của đường thẳng $AB$.}{d}
  \y{Đường thẳng $AB$ có phương trình $5x-8y+10=0$.}{d}
  \y{Điểm $(14;10)$ thuộc đường thẳng $AB$.}{d}
  \y{Đường thẳng $AB$ đi qua gốc toạ độ.}{s}}

\cauTLN[id=T10-C7-B19-D45-TH-063,muc=TH,loigiai={Phương trình $x+6y+30=0$, chia cho $1$: $b=6$, $c=30$.}]{Đường thẳng đi qua $A(6;-6)$ và $B(0;-5)$ có phương trình dạng $x+by+c=0$. Tính $b+c$.}{36}

\cauTLN[id=T10-C7-B19-D45-TH-064,muc=TH,loigiai={$d\colon x-2y-16=0$.}]{Đường thẳng $d$ đi qua $M(6;-5)$ và vuông góc với $2x+y+2=0$ có phương trình dạng $x+by+c=0$. Tính $b+c$.}{-18}

\cauTLN[id=T10-C7-B19-D45-TH-065,muc=TH,loigiai={$d\colon x=0$.}]{Đường thẳng $d$ đi qua $M(0;-1)$ và vuông góc với $6y+1=0$ có phương trình dạng $x+by+c=0$. Tính $b+c$.}{0}

\cauTLN[id=T10-C7-B19-D45-TH-066,muc=TH,loigiai={Phương trình $6x+5y+1=0$, chia cho $6$: $b=\dfrac{5}{6}$, $c=\dfrac{1}{6}$.}]{Đường thẳng đi qua $A(-1;1)$ và $B(4;-5)$ có phương trình dạng $x+by+c=0$. Tính $b+c$.}{1}

\cauTLN[id=T10-C7-B19-D45-TH-067,muc=TH,loigiai={$d\colon x+y=0$.}]{Đường thẳng $d$ đi qua $M(-2;2)$ và vuông góc với $2x-2y+1=0$ có phương trình dạng $x+by+c=0$. Tính $b+c$.}{1}

\cauTLN[id=T10-C7-B19-D45-TH-068,muc=TH,loigiai={$d\colon x-5y-23=0$.}]{Đường thẳng $d$ đi qua $M(3;-4)$ và vuông góc với $5x+y+7=0$ có phương trình dạng $x+by+c=0$. Tính $b+c$.}{-28}

\cauTLN[id=T10-C7-B19-D45-TH-069,muc=TH,loigiai={$d\colon 3x-2y+8=0$.}]{Đường thẳng $d$ đi qua $M(-2;1)$ và vuông góc với $2x+3y-4=0$ có phương trình dạng $x+by+c=0$. Tính $b+c$.}{2}

\cauTLN[id=T10-C7-B19-D45-TH-070,muc=TH,loigiai={Phương trình $x+2=0$, chia cho $1$: $b=0$, $c=2$.}]{Đường thẳng đi qua $A(-2;2)$ và $B(-2;-2)$ có phương trình dạng $x+by+c=0$. Tính $b+c$.}{2}

\cauTN[id=T10-C7-B19-D45-VD-071,muc=VD,loigiai={Trung điểm $BC$: $M(-1;-\dfrac{9}{2})$; đường thẳng $AM\colon 13x-2y+4=0$.}]{Cho tam giác $ABC$ với $A(0;2)$, $B(-6;-6)$, $C(4;-3)$. Phương trình đường trung tuyến kẻ từ $A$ là}
  {$13x-2y+4=0$.}
  {$13x-2y-4=0$.}{$2x-13y-4=0$.}{$13x+2y+4=0$.}

\cauTN[id=T10-C7-B19-D45-VD-072,muc=VD,loigiai={$\overrightarrow{AB}=(7;-7)$, phương trình $x+y=0$.}]{Trên bản đồ một khu công nghiệp (gắn hệ trục $Oxy$, đơn vị $100$ m), nhà máy nước đặt tại $A(-1;1)$ và khu dân cư tại $B(6;-6)$. Người ta làm một đường ống dẫn nước thẳng nối $A$ với $B$. Phương trình đường thẳng chứa đường ống là}
  {$x+y=0$.}
  {$x+y+1=0$.}{$x-y=0$.}{$x+y+2=0$.}

\cauTN[id=T10-C7-B19-D45-VD-073,muc=VD,loigiai={Đi qua trung điểm $I(-\dfrac{11}{2};1)$, vectơ pháp tuyến $\overrightarrow{AB}=(-1;10)$: $2x-20y+31=0$.}]{Phương trình đường trung trực của đoạn thẳng $AB$ với $A(-5;-4)$, $B(-6;6)$ là}
  {$2x-20y+31=0$.}
  {$2x-20y-31=0$.}{$20x-2y-31=0$.}{$2x+20y+31=0$.}

\cauTN[id=T10-C7-B19-D45-VD-074,muc=VD,loigiai={Trung điểm $BC$: $M(-\dfrac{1}{2};5)$; đường thẳng $AM\colon 4x-9y+47=0$.}]{Cho tam giác $ABC$ với $A(-5;3)$, $B(-5;4)$, $C(4;6)$. Phương trình đường trung tuyến kẻ từ $A$ là}
  {$4x-9y+47=0$.}
  {$4x-9y-47=0$.}{$9x-4y-47=0$.}{$4x+9y+47=0$.}

\cauTN[id=T10-C7-B19-D45-VD-075,muc=VD,loigiai={$\overrightarrow{AB}=(13;-6)$, phương trình $6x+13y-3=0$.}]{Trên bản đồ một khu công nghiệp (gắn hệ trục $Oxy$, đơn vị $100$ m), nhà máy nước đặt tại $A(-6;3)$ và khu dân cư tại $B(7;-3)$. Người ta làm một đường ống dẫn nước thẳng nối $A$ với $B$. Phương trình đường thẳng chứa đường ống là}
  {$6x+13y-3=0$.}
  {$6x+13y+3=0$.}{$13x+6y-3=0$.}{$6x-13y-3=0$.}

\cauTN[id=T10-C7-B19-D45-VD-076,muc=VD,loigiai={Đi qua trung điểm $I(\dfrac{7}{2};-\dfrac{1}{2})$, vectơ pháp tuyến $\overrightarrow{AB}=(-1;-3)$: $x+3y-2=0$.}]{Phương trình đường trung trực của đoạn thẳng $AB$ với $A(4;1)$, $B(3;-2)$ là}
  {$x+3y-2=0$.}
  {$x+3y+2=0$.}{$3x+y-2=0$.}{$x-3y-2=0$.}

\cauTN[id=T10-C7-B19-D45-VD-077,muc=VD,loigiai={Trung điểm $BC$: $M(-\dfrac{1}{2};\dfrac{1}{2})$; đường thẳng $AM\colon x+y=0$.}]{Cho tam giác $ABC$ với $A(-3;3)$, $B(0;-1)$, $C(-1;2)$. Phương trình đường trung tuyến kẻ từ $A$ là}
  {$x+y=0$.}
  {$x+y+1=0$.}{$x-y=0$.}{$x+y+2=0$.}

\cauTN[id=T10-C7-B19-D45-VD-078,muc=VD,loigiai={$\overrightarrow{AB}=(10;-6)$, phương trình $3x+5y=0$.}]{Trên bản đồ một khu công nghiệp (gắn hệ trục $Oxy$, đơn vị $100$ m), nhà máy nước đặt tại $A(-5;3)$ và khu dân cư tại $B(5;-3)$. Người ta làm một đường ống dẫn nước thẳng nối $A$ với $B$. Phương trình đường thẳng chứa đường ống là}
  {$3x+5y=0$.}
  {$3x+5y+1=0$.}{$5x+3y=0$.}{$3x-5y=0$.}

\cauTN[id=T10-C7-B19-D45-VD-079,muc=VD,loigiai={Đi qua trung điểm $I(\dfrac{9}{2};-2)$, vectơ pháp tuyến $\overrightarrow{AB}=(-1;4)$: $2x-8y-25=0$.}]{Phương trình đường trung trực của đoạn thẳng $AB$ với $A(5;-4)$, $B(4;0)$ là}
  {$2x-8y-25=0$.}
  {$2x-8y+25=0$.}{$8x-2y+25=0$.}{$2x+8y-25=0$.}

\cauTN[id=T10-C7-B19-D45-VD-080,muc=VD,loigiai={Trung điểm $BC$: $M(-5;\dfrac{9}{2})$; đường thẳng $AM\colon 21x+2y+96=0$.}]{Cho tam giác $ABC$ với $A(-4;-6)$, $B(-6;3)$, $C(-4;6)$. Phương trình đường trung tuyến kẻ từ $A$ là}
  {$21x+2y+96=0$.}
  {$21x+2y-96=0$.}{$2x+21y+96=0$.}{$21x-2y+96=0$.}

\cauTN[id=T10-C7-B19-D45-VD-081,muc=VD,loigiai={$\overrightarrow{AB}=(12;-5)$, phương trình $5x+12y-23=0$.}]{Trên bản đồ một khu công nghiệp (gắn hệ trục $Oxy$, đơn vị $100$ m), nhà máy nước đặt tại $A(-5;4)$ và khu dân cư tại $B(7;-1)$. Người ta làm một đường ống dẫn nước thẳng nối $A$ với $B$. Phương trình đường thẳng chứa đường ống là}
  {$5x+12y-23=0$.}
  {$5x+12y+23=0$.}{$12x+5y-23=0$.}{$5x-12y-23=0$.}

\cauTN[id=T10-C7-B19-D45-VD-082,muc=VD,loigiai={Đi qua trung điểm $I(-3;\dfrac{7}{2})$, vectơ pháp tuyến $\overrightarrow{AB}=(-4;1)$: $8x-2y+31=0$.}]{Phương trình đường trung trực của đoạn thẳng $AB$ với $A(-1;3)$, $B(-5;4)$ là}
  {$8x-2y+31=0$.}
  {$8x-2y-31=0$.}{$2x-8y-31=0$.}{$8x+2y+31=0$.}

\cauDS[id=T10-C7-B19-D45-VD-083,muc=VD,loigiai={$BC\colon 7x+11y+24=0$; đường cao từ $A$: $11x-7y-31=0$. a) Đúng. b) Đúng. c) Đúng. d) Sai.}]{Cho tam giác $ABC$ với $A(-1;-6)$, $B(6;-6)$, $C(-5;1)$.}{%
  \y{Đường thẳng $BC$ có phương trình $7x+11y+24=0$.}{d}
  \y{Đường cao kẻ từ $A$ có vectơ pháp tuyến $\overrightarrow{BC}=(-11;7)$.}{d}
  \y{Đường cao kẻ từ $A$ có phương trình $11x-7y-31=0$.}{d}
  \y{Đường cao kẻ từ $A$ song song với $BC$.}{s}}

\cauDS[id=T10-C7-B19-D45-VD-084,muc=VD,loigiai={$BC\colon 6x+7y-10=0$; đường cao từ $A$: $7x-6y-72=0$. a) Đúng. b) Đúng. c) Đúng. d) Sai.}]{Cho tam giác $ABC$ với $A(6;-5)$, $B(-3;4)$, $C(4;-2)$.}{%
  \y{Đường thẳng $BC$ có phương trình $6x+7y-10=0$.}{d}
  \y{Đường cao kẻ từ $A$ có vectơ pháp tuyến $\overrightarrow{BC}=(7;-6)$.}{d}
  \y{Đường cao kẻ từ $A$ có phương trình $7x-6y-72=0$.}{d}
  \y{Đường cao kẻ từ $A$ song song với $BC$.}{s}}

\cauDS[id=T10-C7-B19-D45-VD-085,muc=VD,loigiai={$BC\colon 9x-y+3=0$; đường cao từ $A$: $x+9y-40=0$. a) Đúng. b) Đúng. c) Đúng. d) Sai.}]{Cho tam giác $ABC$ với $A(4;4)$, $B(-1;-6)$, $C(0;3)$.}{%
  \y{Đường thẳng $BC$ có phương trình $9x-y+3=0$.}{d}
  \y{Đường cao kẻ từ $A$ có vectơ pháp tuyến $\overrightarrow{BC}=(1;9)$.}{d}
  \y{Đường cao kẻ từ $A$ có phương trình $x+9y-40=0$.}{d}
  \y{Đường cao kẻ từ $A$ song song với $BC$.}{s}}

\cauDS[id=T10-C7-B19-D45-VD-086,muc=VD,loigiai={$BC\colon 3x+y+4=0$; đường cao từ $A$: $x-3y-18=0$. a) Đúng. b) Đúng. c) Đúng. d) Sai.}]{Cho tam giác $ABC$ với $A(3;-5)$, $B(-2;2)$, $C(0;-4)$.}{%
  \y{Đường thẳng $BC$ có phương trình $3x+y+4=0$.}{d}
  \y{Đường cao kẻ từ $A$ có vectơ pháp tuyến $\overrightarrow{BC}=(2;-6)$.}{d}
  \y{Đường cao kẻ từ $A$ có phương trình $x-3y-18=0$.}{d}
  \y{Đường cao kẻ từ $A$ song song với $BC$.}{s}}

\cauTLN[id=T10-C7-B19-D45-VD-087,muc=VD,loigiai={$AB\colon 3x-y+12=0$ cắt $Ox$ tại $(-4;0)$, $Oy$ tại $(0;12)$; diện tích $\dfrac12\cdot4\cdot12=24$.}]{Đường thẳng $AB$ với $A(-6;-6)$, $B(-4;0)$ cắt hai trục toạ độ tạo thành một tam giác. Tính diện tích tam giác đó.}{24}

\cauTLN[id=T10-C7-B19-D45-VD-088,muc=VD,loigiai={$d\colon 5x+2y+8=0$.}]{Đường thẳng $d$ đi qua $M(-2;1)$ và vuông góc với $2x-5y+12=0$ có phương trình dạng $x+by+c=0$. Tính $b+c$.}{2}

\cauTLN[id=T10-C7-B19-D45-VD-089,muc=VD,loigiai={$d\colon 4x+5y+23=0$.}]{Đường thẳng $d$ đi qua $M(-2;-3)$ và vuông góc với $5x-4y-2=0$ có phương trình dạng $x+by+c=0$. Tính $b+c$.}{7}

\cauTLN[id=T10-C7-B19-D45-VD-090,muc=VD,loigiai={$d\colon x+y-1=0$.}]{Đường thẳng $d$ đi qua $M(-4;5)$ và vuông góc với $5x-5y-3=0$ có phương trình dạng $x+by+c=0$. Tính $b+c$.}{0}

\cauTN[id=T10-C7-B19-D45-VDC-091,muc=VDC,loigiai={Tam giác $OAB$ vuông cân với $A$, $B$ trên hai tia dương nên $OA=OB=a>0$: $\dfrac xa+\dfrac ya=1$; qua $M$ nên $a=5$.}]{Đường thẳng $d$ đi qua $M(3;2)$ cắt hai tia $Ox$, $Oy$ lần lượt tại $A$, $B$ (khác $O$) sao cho tam giác $OAB$ vuông cân. Phương trình của $d$ là}
  {$x+y-5=0$.}
  {$x-y-1=0$.}{$2x+3y-12=0$.}{$x+y+5=0$.}

\cauTN[id=T10-C7-B19-D45-VDC-092,muc=VDC,loigiai={Trung điểm $BC$: $M(1;-1)$; đường thẳng $AM\colon y+1=0$.}]{Cho tam giác $ABC$ với $A(-2;-1)$, $B(1;-2)$, $C(1;0)$. Phương trình đường trung tuyến kẻ từ $A$ là}
  {$y+1=0$.}
  {$y-1=0$.}{$x+1=0$.}{$y+2=0$.}

\cauTN[id=T10-C7-B19-D45-VDC-093,muc=VDC,loigiai={Tam giác $OAB$ vuông cân với $A$, $B$ trên hai tia dương nên $OA=OB=a>0$: $\dfrac xa+\dfrac ya=1$; qua $M$ nên $a=6$.}]{Đường thẳng $d$ đi qua $M(4;2)$ cắt hai tia $Ox$, $Oy$ lần lượt tại $A$, $B$ (khác $O$) sao cho tam giác $OAB$ vuông cân. Phương trình của $d$ là}
  {$x+y-6=0$.}
  {$x-y-2=0$.}{$x+2y-8=0$.}{$x+y+6=0$.}

\cauTN[id=T10-C7-B19-D45-VDC-094,muc=VDC,loigiai={Trung điểm $BC$: $M(-3;3)$; đường thẳng $AM\colon x+y=0$.}]{Cho tam giác $ABC$ với $A(-5;5)$, $B(0;2)$, $C(-6;4)$. Phương trình đường trung tuyến kẻ từ $A$ là}
  {$x+y=0$.}
  {$x+y+1=0$.}{$x-y=0$.}{$x+y+2=0$.}

\cauTN[id=T10-C7-B19-D45-VDC-095,muc=VDC,loigiai={Tam giác $OAB$ vuông cân với $A$, $B$ trên hai tia dương nên $OA=OB=a>0$: $\dfrac xa+\dfrac ya=1$; qua $M$ nên $a=9$.}]{Đường thẳng $d$ đi qua $M(4;5)$ cắt hai tia $Ox$, $Oy$ lần lượt tại $A$, $B$ (khác $O$) sao cho tam giác $OAB$ vuông cân. Phương trình của $d$ là}
  {$x+y-9=0$.}
  {$x-y+1=0$.}{$5x+4y-40=0$.}{$x+y+9=0$.}

\cauTN[id=T10-C7-B19-D45-VDC-096,muc=VDC,loigiai={Trung điểm $BC$: $M(-\dfrac{3}{2};-\dfrac{9}{2})$; đường thẳng $AM\colon x-3y-12=0$.}]{Cho tam giác $ABC$ với $A(3;-3)$, $B(-6;-3)$, $C(3;-6)$. Phương trình đường trung tuyến kẻ từ $A$ là}
  {$x-3y-12=0$.}
  {$x-3y+12=0$.}{$3x-y+12=0$.}{$x+3y-12=0$.}

\cauDS[id=T10-C7-B19-D45-VDC-097,muc=VDC,loigiai={$BC\colon 9x+y+13=0$; đường cao từ $A$: $x-9y-14=0$. a) Đúng. b) Đúng. c) Sai. d) Sai.}]{Cho tam giác $ABC$ với $A(5;-1)$, $B(-1;-4)$, $C(-2;5)$.}{%
  \y{Đường thẳng $BC$ có phương trình $9x+y+13=0$.}{d}
  \y{Đường cao kẻ từ $A$ có vectơ pháp tuyến $\overrightarrow{BC}=(-1;9)$.}{d}
  \y{Đường cao kẻ từ $A$ có phương trình $x-9y-13=0$.}{s}
  \y{Đường cao kẻ từ $A$ song song với $BC$.}{s}}

\cauDS[id=T10-C7-B19-D45-VDC-098,muc=VDC,loigiai={$BC\colon y-3=0$; đường cao từ $A$: $x+2=0$. a) Đúng. b) Đúng. c) Đúng. d) Sai.}]{Cho tam giác $ABC$ với $A(-2;-2)$, $B(-2;3)$, $C(4;3)$.}{%
  \y{Đường thẳng $BC$ có phương trình $y-3=0$.}{d}
  \y{Đường cao kẻ từ $A$ có vectơ pháp tuyến $\overrightarrow{BC}=(6;0)$.}{d}
  \y{Đường cao kẻ từ $A$ có phương trình $x+2=0$.}{d}
  \y{Đường cao kẻ từ $A$ song song với $BC$.}{s}}

\cauTLN[id=T10-C7-B19-D45-VDC-099,muc=VDC,loigiai={$AB\colon 2x-3y=0$ cắt $Ox$ tại $(0;0)$, $Oy$ tại $(0;0)$; diện tích $\dfrac12\cdot0\cdot0=0$.}]{Đường thẳng $AB$ với $A(-6;-4)$, $B(0;0)$ cắt hai trục toạ độ tạo thành một tam giác. Tính diện tích tam giác đó.}{0}

\cauTLN[id=T10-C7-B19-D45-VDC-100,muc=VDC,loigiai={$AB\colon 10x+y-5=0$ cắt $Ox$ tại $(\dfrac{1}{2};0)$, $Oy$ tại $(0;5)$; diện tích $\dfrac12\cdot\dfrac{1}{2}\cdot5=\dfrac{5}{4}$.}]{Đường thẳng $AB$ với $A(1;-5)$, $B(0;5)$ cắt hai trục toạ độ tạo thành một tam giác. Tính diện tích tam giác đó.}{1,25}
